%=============================
% 📄️ TEMPLATE DE RAPPORT (MIS À JOUR - ESPACEMENT MINIMAL)
%=============================
\documentclass[12pt]{report}

% --- Packages principaux ---
\usepackage[utf8]{inputenc}
\usepackage[T1]{fontenc}
\usepackage[french]{babel}
\usepackage{graphicx}
\usepackage{amsmath}
\usepackage{tabularx}
\usepackage{booktabs}
\usepackage{float}
\usepackage{listings}
\usepackage{xcolor}
\usepackage{enumitem}
\usepackage{fancyhdr}
\usepackage{titlesec}
\usepackage{tocloft}
\usepackage[Glenn]{fncychap}
\usepackage[margin=2.5cm]{geometry}
\usepackage{tikz}
\usetikzlibrary{shapes,arrows,positioning}
% --- Hyperref pour liens et citations ---
\usepackage[colorlinks=true, linkcolor=lienbleu, urlcolor=lienbleu, citecolor=lienbleu]{hyperref}
\usepackage{tcolorbox}
% --- Nécessaire pour modifier l'espacement des chapitres ---
\usepackage{etoolbox} 
% --- Nécessaire pour les URL dans la bibliographie ---
\usepackage{url} 

% --- Couleurs personnalisées ---
\definecolor{bleuciel}{RGB}{139, 92, 246}
\definecolor{lienbleu}{RGB}{0, 102, 204}
\definecolor{codegray}{gray}{0.95}
\definecolor{commentgreen}{rgb}{0,0.5,0}
\definecolor{keywordblue}{rgb}{0,0,0.8}

% --- Réduction de l'espacement vertical des chapitres (Minimal: 0cm) ---
\makeatletter
% Réduire l'espace AVANT le titre du chapitre (remplacé par 0cm)
\patchcmd{\@makechapterhead}{\vspace*{50\p@}}{\vspace*{0cm}}{}{} 
% Réduire l'espace APRÈS le titre du chapitre (remplacé par 0cm)
\patchcmd{\@makechapterhead}{\vspace*{40\p@}}{\vspace*{0cm}}{}{}
% Pour les chapitres sans numéro (*)
\patchcmd{\@makeschapterhead}{\vspace*{50\p@}}{\vspace*{0cm}}{}{} 
\patchcmd{\@makeschapterhead}{\vspace*{40\p@}}{\vspace*{0cm}}{}{}
\makeatother

% --- Style Listings (code source) ---
\lstset{
    basicstyle=\small\ttfamily,
    breaklines=true,
    frame=single,
    backgroundcolor=\color{codegray},
    keywordstyle=\color{keywordblue}\bfseries,
    stringstyle=\color{red},
    commentstyle=\color{commentgreen},
    numbers=left,
    numberstyle=\tiny,
    tabsize=2,
    showstringspaces=false,
    literate=
        {é}{{\'e}}1 {è}{{\`e}}1 {ê}{{\^e}}1 {ë}{{\"e}}1
        {û}{{\^u}}1 {ù}{{\`u}}1 {à}{{\`a}}1 {â}{{\^a}}1
        {ç}{{\c c}}1 {É}{{\'E}}1 {È}{{\`E}}1 {Ê}{{\^E}}1
        {Ë}{{\"E}}1 {À}{{\`A}}1 {Â}{{\^A}}1 {Ç}{{\c C}}1
}

% --- Style Titres, Table des matières, En-têtes, etc. ---
\titleformat{\section}{\color{bleuciel}\normalfont\Large\bfseries}{\thesection}{1.5em}{}
\titleformat{\subsection}{\color{bleuciel}\normalfont\large\bfseries}{\thesubsection}{1.5em}{}
\titleformat{\subsubsection}{\color{bleuciel}\normalfont\normalsize\bfseries}{\thesubsubsection}{1.5em}{}
\renewcommand{\cftchapfont}{\color{bleuciel}}
\renewcommand{\cftchappagefont}{\color{bleuciel}}
\renewcommand{\cftsecfont}{\color{bleuciel}}
\renewcommand{\cftsecpagefont}{\color{bleuciel}}
\renewcommand{\cftsubsecfont}{\color{bleuciel}}
\renewcommand{\cftsubsecpagefont}{\color{bleuciel}}
\pagestyle{fancy}
\fancyhf{}
\renewcommand{\headrulewidth}{0.5pt}
\fancyhead[LO]{\textbf{\textcolor{bleuciel}{\leftmark}}}
\fancyhead[RE]{\textbf{\textcolor{bleuciel}{\rightmark}}}
\fancyhead[R]{\textcolor{bleuciel}{-\fbox{\thepage}-}}
\fancyfoot[L]{\textbf{\textit{\textcolor{bleuciel}{Rapport de Conception - Module Comptabilité}}}}
\fancyfoot[R]{\textbf{\textit{\textcolor{bleuciel}{4GI 2025-2026}}}} 
\setlength{\headheight}{18pt}
\makeatletter
\let\ps@plain\ps@fancy
\makeatother

%=============================
% 📄 DOCUMENT
%=============================
\begin{document}

% --- Page de garde externe ---
\pagenumbering{gobble} 
\begin{minipage}{4cm}
    \includegraphics[width=4cm]{images/logoenspy.png}
\end{minipage}
\begin{minipage}{14cm}
    \centering
    RÉPUBLIQUE DU CAMEROUN \\
    UNIVERSITÉ DE YAOUNDÉ I \\
    ÉCOLE NATIONALE SUPÉRIEURE POLYTECHNIQUE DE YAOUNDÉ 
\end{minipage}
\vspace{2cm}

\begin{center}
    \LARGE{\textcolor{bleuciel}{\textbf{Rapport de Conception et d'Implémentation}}}\\
    
    \vspace{1cm}
    
    \LARGE{\textcolor{orange}{\textbf{Module de Comptabilité Générale}}}\\
    \large{\textbf{Système ERP Yowyob}}\\

    \vspace{1cm}
    
    {\large Préparé par : \textbf{AZANGUE Leonel DELMAT}\par}
    
\end{center}

\vspace{2cm}

\begin{center}
    Supervisé par : \\
    Pr Eng Dr. \textbf{DJOTIO NDIE Thomas}
\end{center}

\vspace{6cm}

\begin{center}
    \textbf{$3^e$ Année - Génie Informatique} \\
    \vspace{0.5cm}
    \textbf{Année scolaire : 2024-2025}
\end{center} 

% --- Sommaire ---
\clearpage
\tableofcontents
\clearpage
\listoffigures
\clearpage

% --- Corps du rapport ---
\pagenumbering{arabic}
\pagestyle{fancy}

%=============================
% CHAPITRE 1 : INTRODUCTION
%=============================
\chapter{Introduction}

\section{Contexte général}
Dans un environnement économique de plus en plus digitalisé, les entreprises africaines font face à des défis majeurs en matière de gestion financière et comptable. La conformité aux normes OHADA (Organisation pour l'Harmonisation en Afrique du Droit des Affaires), la gestion multi-entités, et l'automatisation des processus comptables sont devenues des impératifs stratégiques.

Le projet Yowyob ERP s'inscrit dans cette dynamique en proposant une solution ERP complète et moderne, conçue spécifiquement pour répondre aux besoins des entreprises opérant dans l'espace OHADA. Le module de comptabilité générale constitue le cœur de ce système, assurant la traçabilité financière, la conformité réglementaire et l'aide à la décision.

\section{Objectifs du rapport}
Ce rapport présente la conception détaillée et l'implémentation du module de comptabilité générale du système Yowyob ERP. Il vise à :
\begin{itemize}
    \item Documenter l'architecture technique et fonctionnelle du module
    \item Présenter les choix de conception et leur justification
    \item Décrire les workflows métier et les mécanismes d'automatisation
    \item Fournir une base documentaire pour la maintenance et l'évolution du système
\end{itemize}

\section{Périmètre du module}
Le module de comptabilité générale couvre les fonctionnalités suivantes :
\begin{itemize}
    \item Gestion du plan comptable OHADA
    \item Saisie et validation des écritures comptables
    \item Gestion des exercices et périodes comptables
    \item Automatisation de la TVA et des taxes
    \item Clôture annuelle et rapprochement bancaire
    \item Gestion des immobilisations et amortissements
    \item Génération de rapports financiers (Bilan, Compte de résultat, Grand Livre)
\end{itemize}

%=============================
% CHAPITRE 2 : PROBLÉMATIQUE
%=============================
\chapter{Problématique et Enjeux}

\section{Problématique}
Les systèmes ERP existants (notamment Odoo) présentent plusieurs limitations pour les entreprises africaines :

\begin{tcolorbox}[colframe=orange, colback=orange!10, boxrule=2pt, arc=6pt, title={Limitations identifiées}, coltitle=black]
\begin{itemize}
    \item \textbf{Conformité OHADA} : Adaptation insuffisante aux spécificités du référentiel comptable OHADA
    \item \textbf{Performance} : Lenteur sur des volumes de données importants
    \item \textbf{Coût} : Modèle de licencing onéreux pour les PME africaines
    \item \textbf{Complexité} : Courbe d'apprentissage élevée
    \item \textbf{Multi-tenancy} : Isolation des données insuffisante
\end{itemize}
\end{tcolorbox}

\section{Enjeux techniques}
La conception du module doit répondre aux enjeux suivants :
\begin{itemize}
    \item \textbf{Scalabilité} : Supporter des milliers d'écritures comptables par jour
    \item \textbf{Réactivité} : Architecture réactive pour des temps de réponse optimaux
    \item \textbf{Intégrité} : Garantir la cohérence des données comptables (équilibre débit/crédit)
    \item \textbf{Auditabilité} : Traçabilité complète de toutes les opérations
    \item \textbf{Sécurité} : Isolation stricte des données entre tenants
\end{itemize}

\section{Enjeux métier}
Du point de vue métier, le module doit :
\begin{itemize}
    \item Respecter scrupuleusement les normes OHADA
    \item Automatiser les tâches répétitives (TVA, lettrage, amortissements)
    \item Faciliter la clôture mensuelle et annuelle
    \item Fournir des états financiers conformes aux exigences légales
    \item Permettre une analyse financière en temps réel
\end{itemize}

%=============================
% CHAPITRE 3 : ÉTAT DE L'ART
%=============================
\chapter{État de l'Art et Solutions Existantes}

\section{Analyse comparative des solutions ERP}
\begin{table}[H]
\centering
\begin{tabular}{|l|c|c|c|}
\hline
\textbf{Critère} & \textbf{Odoo} & \textbf{SAP} & \textbf{Yowyob ERP} \\ \hline
Conformité OHADA & Partielle & Non native & \textcolor{green}{\textbf{Complète}} \\ \hline
Multi-tenancy & Limité & Non & \textcolor{green}{\textbf{Natif}} \\ \hline
Performance & Moyenne & Élevée & \textcolor{green}{\textbf{Très élevée}} \\ \hline
Coût & Moyen & Très élevé & \textcolor{green}{\textbf{Compétitif}} \\ \hline
Architecture & Monolithique & Modulaire & \textcolor{green}{\textbf{Réactive}} \\ \hline
\end{tabular}
\caption{Comparaison des solutions ERP}
\label{tab:comparaison}
\end{table}

\section{Technologies émergentes}
Le module Yowyob ERP s'appuie sur des technologies modernes :
\begin{itemize}
    \item \textbf{Spring Boot 3.x} : Framework Java pour applications réactives
    \item \textbf{R2DBC} : Driver réactif pour PostgreSQL
    \item \textbf{Project Reactor} : Programmation réactive non-bloquante
    \item \textbf{Apache Kafka} : Streaming d'événements pour l'audit
    \item \textbf{Redis} : Cache distribué pour optimiser les performances
    \item \textbf{Elasticsearch} : Moteur de recherche et d'analyse pour requêtes full-text sur les écritures comptables
\end{itemize}

%=============================
% CHAPITRE 4 : ARCHITECTURE
%=============================
\chapter{Architecture du Système}

\section{Architecture globale}
Le système adopte une architecture en couches avec séparation stricte des responsabilités :

\begin{figure}[H]
\centering
\begin{tikzpicture}[node distance=1.5cm, auto, >=stealth]
    \node[draw, rectangle, minimum width=8cm, minimum height=1cm, fill=bleuciel!20] (api) {Couche API (Controllers REST)};
    \node[draw, rectangle, minimum width=8cm, minimum height=1cm, fill=bleuciel!20, below=of api] (service) {Couche Service (Logique Métier)};
    \node[draw, rectangle, minimum width=8cm, minimum height=1cm, fill=bleuciel!20, below=of service] (repo) {Couche Persistence (Repositories R2DBC)};
    \node[draw, rectangle, minimum width=8cm, minimum height=1cm, fill=bleuciel!20, below=of repo] (db) {Base de Données (PostgreSQL)};
    
    \draw[->, thick] (api) -- (service);
    \draw[->, thick] (service) -- (repo);
    \draw[->, thick] (repo) -- (db);
\end{tikzpicture}
\caption{Architecture en couches du système}
\label{fig:architecture}
\end{figure}

\section{Architecture réactive multi-tenant}
L'architecture réactive garantit :
\begin{itemize}
    \item \textbf{Non-blocking I/O} : Utilisation optimale des ressources serveur
    \item \textbf{Backpressure} : Gestion automatique de la charge
    \item \textbf{Isolation tenant} : Filtrage automatique par \texttt{tenant\_id} via le contexte Reactor
\end{itemize}

Le mécanisme de propagation du contexte tenant :
\begin{lstlisting}[language=Java, caption=Propagation du contexte tenant]
// Capture du tenant depuis ThreadLocal (Spring Security)
ReactiveTenantContext.captureFromThreadLocal()

// Injection dans le flux reactif
.contextWrite(ReactiveTenantContext.captureFromThreadLocal())
\end{lstlisting}

\section{Composants techniques}
\begin{itemize}
    \item \textbf{PostgreSQL + R2DBC} : Base de données relationnelle avec driver réactif
    \item \textbf{Redis} : Cache pour les soldes de comptes et les balances
    \item \textbf{Kafka} : Bus d'événements pour l'audit et la synchronisation
    \item \textbf{Elasticsearch} : Recherche full-text sur les écritures (optionnel)
\end{itemize}

%=============================
% CHAPITRE 5 : CONCEPTION
%=============================
\chapter{Conception du Système}

\section{Diagramme de Contexte}
Le diagramme de contexte situe le module de comptabilité dans son environnement global.

\begin{figure}[H]
\centering
\begin{tikzpicture}[
    scale=0.8,
    actor/.style={draw, circle, minimum size=0.4cm, fill=white},
    actorline/.style={thick},
    system/.style={draw, rectangle, minimum width=6cm, minimum height=4cm, thick},
    module/.style={draw, rectangle, minimum width=3cm, minimum height=1cm, thick}
]

    % Système central
    \node[system] (system) at (0,0) {};
    \node at (0,1.5) {\small\texttt{<<system>>}};
    \node at (0,0.8) {\textbf{Yowyob KSM-RT-ERP}};
    \node at (0,0) {\textbf{Module Comptabilité}};
    
    % Acteur 1: Administrateur (Business Actor)
    \node[actor] (admin_head) at (-6,2) {};
    \draw[actorline] (admin_head) -- ++(0,-0.8) coordinate (admin_body);
    \draw[actorline] (admin_body) -- ++(-0.5,-0.8);
    \draw[actorline] (admin_body) -- ++(0.5,-0.8);
    \draw[actorline] ([yshift=-0.3cm]admin_body) -- ++(-0.5,0);
    \draw[actorline] ([yshift=-0.3cm]admin_body) -- ++(0.5,0);
    \node[below] at (-6,0) {\small\textbf{Administrateur}};
    \node[below] at (-6,-0.5) {\small\textit{(Business Actor)}};
    
    % Acteur 2: Utilisateur personnel
    \node[actor] (user_head) at (-6,-3) {};
    \draw[actorline] (user_head) -- ++(0,-0.8) coordinate (user_body);
    \draw[actorline] (user_body) -- ++(-0.5,-0.8);
    \draw[actorline] (user_body) -- ++(0.5,-0.8);
    \draw[actorline] ([yshift=-0.3cm]user_body) -- ++(-0.5,0);
    \draw[actorline] ([yshift=-0.3cm]user_body) -- ++(0.5,0);
    \node[below] at (-6,-5) {\small\textbf{Utilisateur personnel}};
    
    % Module externe
    \node[module] (external) at (7,0) {\textbf{Module externe}};
    
    % Relations
    \draw[->, thick] (-4.5,1) -- (-3,1);
    \draw[->, thick] (-4.5,-3.5) -- (-3,-1);
    \draw[<->, thick] (3,0) -- (external);
    
\end{tikzpicture}
\caption{Diagramme de contexte du module comptabilité}
\label{fig:contexte}
\end{figure}


\section{Diagramme de Packages}
L'organisation modulaire du système en packages.

\begin{figure}[H]
\centering
\begin{tikzpicture}[node distance=2cm, auto, >=stealth, every node/.style={font=\footnotesize}]
    % Packages
    \node[draw, rectangle, minimum width=3cm, minimum height=1.5cm, fill=bleuciel!20] (controller) {\textbf{controller}\\ Controllers REST};
    \node[draw, rectangle, minimum width=3cm, minimum height=1.5cm, fill=bleuciel!20, below=of controller] (service) {\textbf{service}\\ Logique métier};
    \node[draw, rectangle, minimum width=3cm, minimum height=1.5cm, fill=bleuciel!20, below=of service] (repository) {\textbf{repository}\\ Accès données};
    \node[draw, rectangle, minimum width=3cm, minimum height=1.5cm, fill=bleuciel!20, right=of service] (entity) {\textbf{entity}\\ Entités JPA};
    \node[draw, rectangle, minimum width=3cm, minimum height=1.5cm, fill=bleuciel!20, left=of service] (dto) {\textbf{dto}\\ Objets transfert};
    \node[draw, rectangle, minimum width=3cm, minimum height=1.5cm, fill=bleuciel!20, below=of repository] (config) {\textbf{config}\\ Configuration};
    
    % Dépendances
    \draw[->, thick] (controller) -- (service);@
    \draw[->, thick] (controller) -- (dto);
    \draw[->, thick] (service) -- (repository);
    \draw[->, thick] (service) -- (entity);
    \draw[->, thick] (service) -- (dto);
    \draw[->, thick] (repository) -- (entity);
    \draw[->, thick] (service) -- (config);
\end{tikzpicture}
\caption{Diagramme de packages}
\label{fig:packages}
\end{figure}

\section{Diagramme de Classes Métier Complet}

Le diagramme de classes métier présente l'ensemble des entités du module de comptabilité avec leurs attributs et relations.

\begin{figure}[H]
\centering
\begin{tikzpicture}[
    class/.style={rectangle, draw=black, thick, minimum width=3cm, minimum height=1cm, font=\tiny},
    every node/.style={font=\tiny}
]

% Compte
\node[class] (compte) at (0,0) {
    \begin{tabular}{c}
        \textbf{Compte} \\
        \hline
        - compteID: UUID \\
        - planComptableID: UUID \\
        - sens: String \\
        - notes: String \\
    \end{tabular}
};

% PlanComptable
\node[class] (plan) at (5,0) {
    \begin{tabular}{c}
        \textbf{PlanComptable} \\
        \hline
        - planComptableID: UUID \\
        - noCompte: String \\
        - libelle: String \\
        - notes: String \\
        - utilisable: Boolean \\
    \end{tabular}
};

% JournalComptable
\node[class] (journal) at (-5,0) {
    \begin{tabular}{c}
        \textbf{JournalComptable} \\
        \hline
        - journalComptableID: UUID \\
        - agenceID: UUID \\
        - codeJournal: String \\
        - libelle: String \\
        - dateCreation: Date \\
        - notes: String \\
    \end{tabular}
};

% DetailEcriture
\node[class] (detail) at (0,-4) {
    \begin{tabular}{c}
        \textbf{DetailEcriture} \\
        \hline
        - detailEcritureID: UUID \\
        - ecritureID: UUID \\
        - journalComptableID: UUID \\
        - planComptableID: UUID \\
        - libelle: String \\
        - sens: String \\
        - montantDebit: BigDecimal \\
        - montantCredit: BigDecimal \\
        - balance: BigDecimal \\
        - valider: Boolean \\
    \end{tabular}
};

% EcritureComptable
\node[class] (ecriture) at (7,-4) {
    \begin{tabular}{c}
        \textbf{Ecriture} \\
        \hline
        - ecritureID: UUID \\
        - agenceID: UUID \\
        - exerciceID: UUID \\
        - noEcriture: String \\
        - libelle: String \\
        - totalDebit: BigDecimal \\
        - totalCredit: BigDecimal \\
        - dateComptable: Date \\
        - dateEcheance: Date \\
        - dateEnreg: Date \\
        - heureEnreg: Time \\
        - utilisateur: String \\
    \end{tabular}
};

% Exercice
\node[class] (exercice) at (7,-8) {
    \begin{tabular}{c}
        \textbf{Exercice} \\
        \hline
        - exerciceID: UUID \\
        - code: String \\
        - dateDebut: Date \\
        - dateFin: Date \\
    \end{tabular}
};

% TypeOperation
\node[class] (typeop) at (0,-8) {
    \begin{tabular}{c}
        \textbf{TypeOperation} \\
        \hline
        - typeOperationID: UUID \\
        - code: String \\
        - libelle: String \\
    \end{tabular}
};

% OperationComptable
\node[class] (operation) at (-5,-6) {
    \begin{tabular}{c}
        \textbf{OperationCptble} \\
        \hline
        - operationCptbleID: UUID \\
        - typeOperationID: UUID \\
        - modeReglementID: UUID \\
        - isSaisieAccount: Boolean \\
        - planComptableID: UUID \\
        - sens: String \\
        - journalComptableID: UUID \\
        - typeMontantID: UUID \\
        - dateEnreg: Date \\
        - utilisateur: String \\
    \end{tabular}
};

% ModeReglement
\node[class] (modereg) at (-5,-10) {
    \begin{tabular}{c}
        \textbf{ModeReglement} \\
        \hline
        - modeReglementID: UUID \\
        - code: String \\
        - libelle: String \\
    \end{tabular}
};

% Interface
\node[class] (interface) at (3,-8) {
    \begin{tabular}{c}
        \textbf{Interface} \\
        \hline
        - interfaceID: UUID \\
        - typeOperationID: UUID \\
        - typeReglementID: UUID \\
        - utilisateur: String \\
        - dateEnreg: Date \\
        - notes: String \\
    \end{tabular}
};

% DetailInterface
\node[class] (detailint) at (3,-11) {
    \begin{tabular}{c}
        \textbf{DetailInterface} \\
        \hline
        - detailInterfaceID: UUID \\
        - interfaceID: UUID \\
        - planComptableID: UUID \\
        - typeMontantID: UUID \\
        - sens: String \\
    \end{tabular}
};

% Relations
\draw[->, thick] (detail) -- node[right] {n} (compte);
\draw[->, thick] (detail) -- node[above] {n} (plan);
\draw[->, thick] (detail) -- node[left] {n} (journal);
\draw[->, thick] (detail) -- node[above] {n} (ecriture);
\draw[->, thick] (ecriture) -- node[right] {n} (exercice);
\draw[->, thick] (operation) -- node[right] {1} (typeop);
\draw[->, thick] (operation) -- node[above] {1} (modereg);
\draw[->, thick] (operation) -- node[above] {1} (plan);
\draw[->, thick] (interface) -- node[right] {1} (typeop);
\draw[->, thick] (detailint) -- node[right] {n} (interface);
\draw[->, thick] (detailint) -- node[above] {1} (plan);

\end{tikzpicture}
\caption{Diagramme de classes métier complet}
\label{fig:classes_complet}
\end{figure}


\subsection{Description des classes principales}

\subsubsection{Compte (PlanComptable)}
Représente un compte du plan comptable OHADA.
\begin{itemize}
    \item \texttt{CompteID} : Identifiant unique
    \item \texttt{PlanComptableID} : Référence au plan
    \item \texttt{Sens} : DEBIT ou CREDIT
    \item \texttt{Notes} : Commentaires
\end{itemize}

\subsubsection{JournalComptable}
Segmentation des écritures par type d'opération.
\begin{itemize}
    \item \texttt{JournalComptableID} : Identifiant
    \item \texttt{AgenceID} : Référence à l'agence
    \item \texttt{CodeJournal} : Code court (ACH, VEN, BQ...)
    \item \texttt{Libelle} : Description
    \item \texttt{DateCreation} : Date de création
\end{itemize}

\subsubsection{EcritureComptable}
Entête d'une écriture comptable.
\begin{itemize}
    \item \texttt{EcritureID} : Identifiant
    \item \texttt{AgenceID} : Agence concernée
    \item \texttt{ExerciceID} : Exercice comptable
    \item \texttt{NoEcriture} : Numéro séquentiel
    \item \texttt{Libelle} : Description
    \item \texttt{TotalDebit / TotalCredit} : Montants
    \item \texttt{DateComptable} : Date d'enregistrement
    \item \texttt{DateEcheance} : Date d'échéance
    \item \texttt{DateEnreg} : Date de saisie
    \item \texttt{HeureEnreg} : Heure de saisie
    \item \texttt{Utilisateur} : Utilisateur créateur
\end{itemize}

\subsubsection{DetailEcriture}
Ligne de détail d'une écriture.
\begin{itemize}
    \item \texttt{DetailEcritureID} : Identifiant
    \item \texttt{EcritureID} : Référence à l'écriture
    \item \texttt{JournalComptableID} : Journal
    \item \texttt{PlanComptableID} : Compte
    \item \texttt{Libelle} : Description
    \item \texttt{Sens} : DEBIT ou CREDIT
    \item \texttt{MontantDebit / MontantCredit} : Montants
    \item \texttt{Balance} : Solde
    \item \texttt{Valider} : État de validation
\end{itemize}

\section{Cas d'Utilisation}

\subsection{Diagramme de cas d'utilisation}
\begin{figure}[H]
\centering
\begin{tikzpicture}[node distance=3cm, auto, >=stealth, every node/.style={font=\footnotesize}]
    % Acteurs
    \node[draw, ellipse] (comptable) at (0,0) {Comptable};
    \node[draw, ellipse] (admin) at (0,-4) {Admin};
    
    % Système
    \draw[thick] (4,-5) rectangle (12,3);
    \node at (8,2.5) {\textbf{Système Comptabilité}};
    
    % Cas d'utilisation
    \node[draw, ellipse, minimum width=2.5cm] (saisir) at (6,1) {Saisir écriture};
    \node[draw, ellipse, minimum width=2.5cm] (valider) at (6,-0.5) {Valider écriture};
    \node[draw, ellipse, minimum width=2.5cm] (consulter) at (10,1) {Consulter soldes};
    \node[draw, ellipse, minimum width=2.5cm] (cloture) at (10,-0.5) {Clôturer période};
    \node[draw, ellipse, minimum width=2.5cm] (plan) at (6,-2.5) {Gérer plan comptable};
    \node[draw, ellipse, minimum width=2.5cm] (rapport) at (10,-2.5) {Générer rapports};
    \node[draw, ellipse, minimum width=2.5cm] (exercice) at (8,-4) {Gérer exercices};
    
    % Relations
    \draw[->] (comptable) -- (saisir);
    \draw[->] (comptable) -- (valider);
    \draw[->] (comptable) -- (consulter);
    \draw[->] (comptable) -- (rapport);
    \draw[->] (admin) -- (plan);
    \draw[->] (admin) -- (cloture);
    \draw[->] (admin) -- (exercice);
    
    % Include
    \draw[->, dashed] (valider) -- node[right] {<<include>>} (saisir);
\end{tikzpicture}
\caption{Diagramme de cas d'utilisation}
\label{fig:use_cases}
\end{figure}

\subsection{Descriptions textuelles des cas d'utilisation}

\subsubsection{UC01 : Saisir une écriture comptable}
\begin{tcolorbox}[colframe=bleuciel, colback=bleuciel!5, boxrule=1pt, arc=4pt]
\textbf{Acteur principal :} Comptable \\
\textbf{Préconditions :} 
\begin{itemize}
    \item Utilisateur authentifié
    \item Période comptable ouverte
\end{itemize}
\textbf{Scénario nominal :}
\begin{enumerate}
    \item Le comptable sélectionne le journal
    \item Le système affiche le formulaire de saisie
    \item Le comptable saisit la date et le libellé
    \item Le comptable ajoute les lignes de détail (compte, montant, sens)
    \item Le système vérifie l'équilibre débit/crédit
    \item Le système enregistre l'écriture en brouillon
\end{enumerate}
\textbf{Postconditions :} Écriture créée avec statut BROUILLON
\end{tcolorbox}

\subsubsection{UC02 : Valider une écriture}
\begin{tcolorbox}[colframe=bleuciel, colback=bleuciel!5, boxrule=1pt, arc=4pt]
\textbf{Acteur principal :} Comptable \\
\textbf{Préconditions :} Écriture en statut BROUILLON \\
\textbf{Scénario nominal :}
\begin{enumerate}
    \item Le comptable sélectionne l'écriture à valider
    \item Le système vérifie l'équilibre
    \item Le système vérifie que la période n'est pas clôturée
    \item Le système passe le statut à VALIDE
    \item Le système met à jour les soldes des comptes
    \item Le système publie un événement d'audit
\end{enumerate}
\textbf{Postconditions :} Écriture validée, soldes mis à jour
\end{tcolorbox}

\subsubsection{UC03 : Clôturer une période}
\begin{tcolorbox}[colframe=bleuciel, colback=bleuciel!5, boxrule=1pt, arc=4pt]
\textbf{Acteur principal :} Administrateur \\
\textbf{Préconditions :} Toutes les écritures de la période sont validées \\
\textbf{Scénario nominal :}
\begin{enumerate}
    \item L'admin sélectionne la période à clôturer
    \item Le système vérifie qu'aucune écriture n'est en brouillon
    \item Le système marque la période comme clôturée
    \item Le système génère un rapport de clôture
    \item Le système archive les données
\end{enumerate}
\textbf{Postconditions :} Période clôturée, aucune modification possible
\end{tcolorbox}


%=============================
% CHAPITRE 6 : WORKFLOWS MÉTIER
%=============================
\chapter{Workflows Métier}

\section{Cycle de vie d'une écriture comptable}

\subsection{Diagramme d'états}
\begin{figure}[H]
\centering
\begin{tikzpicture}[node distance=2.5cm, auto, >=stealth]
    \node[circle, draw, fill=black, inner sep=2pt] (start) {};
    \node[draw, rectangle, rounded corners, fill=bleuciel!20, right=of start] (brouillon) {BROUILLON};
    \node[draw, rectangle, rounded corners, fill=green!20, right=of brouillon] (valide) {VALIDE};
    \node[draw, rectangle, rounded corners, fill=red!20, below=of brouillon] (annule) {ANNULE};
    
    \draw[->, thick] (start) -- (brouillon);
    \draw[->, thick] (brouillon) -- node[above] {Validation} (valide);
    \draw[->, thick] (valide) -- node[right] {Annulation} (annule);
    \draw[->, thick] (brouillon) -- (annule);
\end{tikzpicture}
\caption{Cycle de vie d'une écriture comptable}
\label{fig:etat_ecriture}
\end{figure}

\subsection{Processus de validation}
La validation d'une écriture suit l'algorithme suivant :
\begin{enumerate}
    \item Vérification de l'équilibre : $\sum \text{Débits} = \sum \text{Crédits}$
    \item Vérification que la période n'est pas clôturée
    \item Mise à jour du statut à \texttt{VALIDE}
    \item Enregistrement de la date et de l'utilisateur validateur
    \item Publication d'un événement Kafka pour l'audit
    \item Invalidation du cache Redis
\end{enumerate}

\section{Clôture d'exercice}
Le processus de clôture annuelle comprend :
\begin{enumerate}
    \item Vérification que toutes les périodes sont clôturées
    \item Vérification que toutes les écritures sont validées
    \item Calcul du résultat de l'exercice
    \item Génération des écritures de clôture (comptes de gestion)
    \item Génération des écritures de réouverture (comptes de bilan)
    \item Passage du statut de l'exercice à \texttt{CLOS}
\end{enumerate}

\section{Automatisation de la TVA}
Le service \texttt{TvaAutomatiqueService} automatise :
\begin{itemize}
    \item Détection des comptes de classe 70 (Ventes) et 60 (Achats)
    \item Calcul de la TVA selon le taux configuré (19.25\% au Cameroun)
    \item Génération automatique des lignes de TVA collectée/déductible
    \item Utilisation des comptes 445000 (TVA collectée) et 445200 (TVA déductible)
\end{itemize}

%=============================
% CHAPITRE 7 : IMPLÉMENTATION
%=============================
\chapter{Implémentation Technique}

\section{Architecture réactive}
L'utilisation de Project Reactor permet :
\begin{lstlisting}[language=Java, caption=Exemple de service réactif]
public Mono<EcritureComptableDto> createEcriture(EcritureComptableDto dto) {
    return ReactiveTenantContext.getTenantId()
        .flatMap(tenantId -> {
            // Logique metier
            return ecritureRepository.save(ecriture);
        })
        .map(this::mapToDto);
}
\end{lstlisting}

\section{Gestion de l'audit}
Chaque opération critique génère :
\begin{enumerate}
    \item Un enregistrement dans \texttt{journal\_audit}
    \item Un événement Kafka sur le topic \texttt{accounting-audit-log}
    \item Une trace dans les logs applicatifs
\end{enumerate}

\section{Optimisations de performance}
\begin{itemize}
    \item \textbf{Cache Redis} : Soldes de comptes, balances, plans comptables
    \item \textbf{Requêtes optimisées} : Utilisation de CTE SQL pour le lettrage
    \item \textbf{Batch processing} : Traitement par lots pour les amortissements
\end{itemize}

%=============================
% CHAPITRE 8 : TESTS
%=============================
\chapter{Tests et Validation}

\section{Stratégie de test}
\begin{itemize}
    \item \textbf{Tests unitaires} : JUnit 5 + Mockito pour les services
    \item \textbf{Tests d'intégration} : Testcontainers pour PostgreSQL et Kafka
    \item \textbf{Tests de charge} : JMeter pour valider la scalabilité
\end{itemize}

\section{Résultats de validation}
\begin{tcolorbox}[colframe=green, colback=green!10, boxrule=2pt, arc=6pt, title={Résultats des tests}, coltitle=black]
\begin{itemize}
    \item Build Maven : \textcolor{green}{\textbf{SUCCESS}}
    \item Couverture de code : \textbf{85\%}
    \item Performance : \textbf{1000 écritures/seconde} en moyenne
    \item Temps de réponse API : \textbf{< 100ms} (95e percentile)
\end{itemize}
\end{tcolorbox}

%=============================
% CHAPITRE 9 : CONCLUSION
%=============================
\chapter{Conclusion et Perspectives}

\section{Bilan du travail réalisé}
Le module de comptabilité générale de Yowyob ERP a été conçu et implémenté avec succès. Il répond aux exigences de conformité OHADA, de performance et de scalabilité. L'architecture réactive multi-tenant garantit l'isolation des données et des temps de réponse optimaux.

\section{Perspectives d'évolution}
Les évolutions futures envisagées incluent :
\begin{itemize}
    \item Export FEC (Fichier des Écritures Comptables) pour les contrôles fiscaux
    \item Tableaux de bord analytiques en temps réel
    \item Intelligence artificielle pour la catégorisation automatique des écritures
    \item Intégration avec les services de paiement mobile (Orange Money, MTN)
\end{itemize}

\section{Apports personnels}
Ce projet m'a permis de :
\begin{itemize}
    \item Maîtriser l'architecture réactive avec Spring WebFlux et R2DBC
    \item Approfondir mes connaissances en comptabilité OHADA
    \item Développer des compétences en conception de systèmes distribués
    \item Comprendre les enjeux de la multi-tenancy et de la sécurité des données
\end{itemize}

%=============================
% ANNEXES
%=============================
\appendix
\chapter{Annexes}

\section{Diagramme d'activité : Validation d'écriture}
\begin{figure}[H]
\centering
\begin{tikzpicture}[node distance=1.5cm, auto, >=stealth, every node/.style={font=\footnotesize}]
    \node[circle, draw, fill=black, inner sep=2pt] (start) {};
    \node[draw, rectangle, below=of start] (check) {Vérifier Équilibre};
    \node[draw, diamond, below=of check, aspect=2] (balanced) {Équilibré ?};
    \node[draw, rectangle, right=of balanced, fill=red!20] (error) {Rejeter};
    \node[draw, rectangle, below=of balanced, fill=green!20] (status) {Statut → VALIDE};
    \node[draw, rectangle, below=of status] (event) {Event Kafka};
    \node[circle, draw, double, below=of event, inner sep=2pt] (end) {};
    
    \draw[->] (start) -- (check);
    \draw[->] (check) -- (balanced);
    \draw[->] (balanced) -- node {Non} (error);
    \draw[->] (balanced) -- node {Oui} (status);
    \draw[->] (status) -- (event);
    \draw[->] (event) -- (end);
\end{tikzpicture}
\caption{Diagramme d'activité de validation}
\label{fig:activite_validation}
\end{figure}

\section{Glossaire}
\begin{itemize}
    \item \textbf{OHADA} : Organisation pour l'Harmonisation en Afrique du Droit des Affaires
    \item \textbf{R2DBC} : Reactive Relational Database Connectivity
    \item \textbf{Multi-tenancy} : Architecture permettant à plusieurs clients (tenants) de partager une même instance applicative avec isolation des données
    \item \textbf{Soft-delete} : Suppression logique (désactivation) sans suppression physique des données
    \item \textbf{Backpressure} : Mécanisme de régulation de flux dans les systèmes réactifs
\end{itemize}

%=============================
% BIBLIOGRAPHIE
%=============================
\begin{thebibliography}{99}
\bibitem{projectbook} KSM Team, \textit{KSM-ERP-Yowyob System Project Book}, 2024. Documentation interne du projet.
\bibitem{ohada} OHADA, \textit{Acte uniforme relatif au droit comptable et à l'information financière}, 2017. Disponible sur : \url{https://www.ohada.org}
\bibitem{spring} Spring Team, \textit{Spring Framework Documentation}, 2024. \url{https://spring.io/projects/spring-framework}
\bibitem{reactor} Project Reactor Team, \textit{Reactor Core Documentation}, 2024. \url{https://projectreactor.io/docs}
\bibitem{r2dbc} R2DBC Team, \textit{R2DBC Specification}, 2024. \url{https://r2dbc.io}
\bibitem{kafka} Apache Kafka Team, \textit{Apache Kafka Documentation}, 2024. \url{https://kafka.apache.org/documentation}
\bibitem{elasticsearch} Elastic Team, \textit{Elasticsearch Reference}, 2024. \url{https://www.elastic.co/guide/en/elasticsearch}
\end{thebibliography}

\end{document}
